\documentclass[conference,onecolumn]{IEEEtran}
% =========================================================
% Arc A paper: "The Gate Is Only as Honest as Its Contracts"
% Companion to main.tex (Capability Minimization as a Safety Primitive / RACG).
% Reuses the same preamble conventions and references.bib.
% =========================================================
\usepackage[T1]{fontenc}
\usepackage[utf8]{inputenc}
\usepackage{newtxtext}
\usepackage{cite}
\usepackage{amsmath,amssymb,amsfonts}
\usepackage{amsthm}
\newtheorem{proposition}{Proposition}
\usepackage{newtxmath}
\usepackage{algorithm}
\usepackage{algorithmic}
\usepackage{graphicx}
\usepackage{textcomp}
\usepackage{xcolor}
\usepackage{float}
\usepackage{placeins}
\usepackage{array}
\usepackage{booktabs}
\usepackage{multirow}
\usepackage{tabularx}
\usepackage{ragged2e}
\usepackage{url}
\newcolumntype{L}[1]{>{\RaggedRight\arraybackslash}p{#1}}
\newcolumntype{Y}{>{\RaggedRight\arraybackslash}X}
\renewcommand{\arraystretch}{1.15}
\setlength{\tabcolsep}{4pt}
\setlength{\emergencystretch}{2em}

% shorthand
\newcommand{\racg}{\textsc{RACG}}
\newcommand{\cg}{\textsc{ContractGuard}}
\newcommand{\isr}{\mathrm{ISR}}
\newcommand{\Vt}{V_t}
\newcommand{\code}[1]{\texttt{#1}}

\begin{document}

\title{The Gate Is Only as Honest as Its Contracts:\\
ContractGuard for the Contract Layer of Risk-Aware Causal Gating}

\author{\IEEEauthorblockN{Laxmipriya Ganesh Iyer}
\IEEEauthorblockA{\textit{Independent Researcher} \\
United States of America \\
iyer.la@northeastern.edu}
\and
\IEEEauthorblockN{Rahul Suresh Babu}
\IEEEauthorblockA{\textit{Independent Researcher} \\
United States of America \\
rahulsb@bu.edu}
}

\maketitle

\begin{abstract}
Risk-Aware Causal Gating (\racg) defends tool-augmented LLM agents against
indirect prompt injection by removing dangerous tools from the agent's visible
action space, so that even a fully injection-compliant agent cannot call a tool
it cannot see. We make three points, which the body develops into five
contributions. \textbf{First}, this structural guarantee
does not eliminate the trust assumption behind safe tool use; it \emph{relocates}
it into the integrity of the tool \emph{contracts}---declared preconditions,
effects, risk, and authorization---that the gate reads, so an attacker who
corrupts a contract (via a poisoned tool description, an untrusted plugin, or a
tool that misreports its behavior) can make the gate mis-decide
\emph{without ever persuading the agent}. \textbf{Second}, forging a tool's
\emph{effects} is strictly more dangerous than tampering with its
\emph{risk label}: \racg{} applies a causal gate before its admissibility gate,
so an off-path tool is never exposed and risk-relabeling alone fails, whereas
effect/precondition forgery routes the dangerous tool onto the causal path and
succeeds. Effect integrity, not the risk label, is the load-bearing assumption.
\textbf{Third}, we introduce \cg, a verifier between the registry and the gate
that layers signed provenance, typed contract attestation, and runtime effect
verification; on a controlled benchmark it restores injection success to zero
against every modeled attack---including an exhaustive white-box adaptive
attacker---without over-rejecting honest contracts, and we show the original
$\lambda$ risk knob cannot substitute for contract integrity.
\end{abstract}

\begin{IEEEkeywords}
LLM agents, agent safety, prompt injection, tool contracts, provenance,
information integrity, least privilege, attack surface, supply-chain security,
function calling, AI safety
\end{IEEEkeywords}

% =========================================================
\section{Introduction}
% =========================================================
Tool-augmented large language model (LLM) agents increasingly hold authority
over high-consequence actions---sending messages, transferring funds, deleting
records~\cite{yao2022react,schick2023toolformer,qin2023toollm}. Risk-Aware
Causal Gating (\racg)~\cite{iyer2026racg} addresses the resulting exposure
problem by treating the visible tool set $\Vt$ as a security control surface and
applying least privilege~\cite{saltzer1975protection}. In one sentence: \racg{}
reads, from each tool's declared \emph{contract}, what the tool does (its
effects), what it needs (its preconditions), how dangerous it is (a risk tier),
and what authorization it requires; it then plans a minimal-cost path to the
user's goal and exposes a high-risk tool \emph{only} when that tool is both
causally necessary for the goal and authorized in the current state, hiding it
otherwise. \racg{} builds on causal minimal tool
filtering~\cite{anon2026cmtf} and on \emph{learned} precondition--effect tool
contracts~\cite{iyer2026contract2tool}, which supply the $R_i,E_i$ annotations
the gate relies on. The headline safety property is \emph{structural}: if the
dangerous tool is not in $\Vt$, then even a worst-case agent that obeys every
injected instruction cannot call it, because the call is not available to
attempt. This makes injection success a property of the \emph{action space}
rather than of the model's willingness to refuse.

We make a simple but consequential observation: this structural guarantee does
not remove the trust assumption behind safe tool use---it \emph{relocates} it.
\racg{} does not \emph{know} that \code{send\_email} is dangerous or that
\code{transfer\_funds} requires confirmation. It \emph{reads} these facts from
the tool's \emph{contract}: a declared risk tier $\rho_i$, a produced-effect set
$E_i$, a required-precondition set $R_i$, and an authorization set
$\alpha_i$~\cite{iyer2026racg}. The gate trusts all of this as ground truth.
Consequently, the safety of the whole system rests on the \emph{integrity of the
contracts}, an assumption the original formulation states only informally as a
provenance condition on authorization variables and otherwise leaves
unexamined. Indeed, the original work explicitly names ``runtime provenance
enforcement (e.g.\ taint tracking or signed authorization facts)'' as future
work; this paper takes up the contract-integrity half of that program.

We therefore threat-model the contract layer itself. Our threat model grants the
attacker influence over contract \emph{declarations and runtime behavior}, but
not over the gate, the agent, or a trusted attestation of the original registry.
This captures realistic supply-chain and content-channel attacks: a
prompt-injected tool description that adds a false effect annotation
(``this tool produces \code{auth\_token}''), a tool registry assembled in part
from untrusted third-party plugins, or a tool whose advertised behavior diverges
from what it actually does at runtime. Crucially, in all of these the attacker
\emph{never tries to persuade the agent}. They corrupt the contract so the gate
makes the wrong decision, and the dangerous tool re-enters $\Vt$; from there even
a compliant agent calls it.

A central finding emerges from instrumenting \racg{} under attack. \racg{}
actually contains \emph{two} gates applied in sequence: a \emph{causal gate} that
asks whether a tool lies on a minimal-cost path to the goal, and an
\emph{admissibility gate} that asks whether a risky frontier tool is authorized.
The causal gate dominates. A tool that does not lie on a minimal path to the goal
is never the frontier tool, and therefore never reaches admissibility. As a
direct corollary, the two ``obvious'' contract attacks---\emph{relabeling} a
dangerous tool as low-risk, or \emph{aliasing away} its authorization
requirement---are, on their own, \textbf{insufficient} to defeat \racg, because
the targeted tool is never causally selected in the first place. This is a
robustness result for the original method, and it sharpens our threat model: the
load-bearing trust is \emph{effect integrity} (the $R_i,E_i$ that determine
causal reachability), not the risk label $\rho_i$. The attack that works must
first defeat the causal gate by forging effects or preconditions so that the
dangerous tool becomes part of the cheapest path to the goal.

Having characterized the vulnerability, we propose a defense. \cg{} is a verifier
interposed between the (possibly poisoned) registry and the gate; it returns a
sanitized registry the gate can safely read. \cg{} composes three mechanisms:
(1)~\emph{signed provenance}, which rejects contract fields not vouched for by a
trusted signer; (2)~\emph{typed contract attestation}, which verifies a
trusted-signed tool's gate-relevant fields (preconditions, effects, risk, and
authorization) against the attestation (subsuming the rule that a tool may only
establish authorization variables it is entitled to);
and (3)~\emph{runtime effect verification}, which
checks a tool's realized output against its declared effects and drops
undeclared effects. We show these form a \emph{strict necessity ladder}: each
mechanism neutralizes a distinct attack and is provably insufficient against the
next, mirroring the original method's hypothesis ladder.

This paper makes five contributions.
\textbf{First}, we formulate contract integrity as a hidden, load-bearing trust
assumption of structural tool gating, and give a threat model for the contract
layer.
\textbf{Second}, we identify and empirically demonstrate the two-gate structure
of \racg{} and the resulting effect-integrity-primary attack taxonomy, including
the negative result that risk-relabeling and authorization-aliasing alone do not
defeat \racg, generalized across the whole registry and quantified by a
per-field ablation.
\textbf{Third}, we introduce \cg{} and show that its three mechanisms---signed
provenance, typed contract attestation, and runtime effect
verification---form a strict necessity ladder restoring injection success to zero
against shortcut forgery, signed over-scoping, runtime divergence, and a compound
attack, without over-rejecting honest contracts, and we give soundness
propositions for each rung.
\textbf{Fourth}, we evaluate against an \emph{exhaustive white-box adaptive}
attacker and show the full stack admits no successful contract perturbation,
while partial deployments are each defeated; this evaluation also uncovered and
fixed a gap in a signer-only formulation of the defense.
\textbf{Fifth}, we show the original safety dial $\lambda$ is orthogonal to
contract integrity: no setting of $\lambda$ removes a contract-forgery attack,
because the attack defeats the causal gate that $\lambda$ does not influence.

% =========================================================
\section{Background and Threat Model}
\label{sec:bg}
% =========================================================

\subsection{Tool Contracts and Risk-Aware Causal Gating}
We adopt the precondition--effect tool-contract formalism of the base
method~\cite{iyer2026racg}, whose contracts are learned and annotated following
Contract2Tool~\cite{iyer2026contract2tool} and filtered causally following
CMTF~\cite{anon2026cmtf}. A tool is a tuple
$t_i = (d_i, R_i, E_i, c_i, \rho_i, \alpha_i)$ with natural-language description
$d_i$, preconditions $R_i$ (state variables required to execute), effects $E_i$
(state variables produced), cost $c_i$, ordinal risk $\rho_i \in
\{\textsc{low},\textsc{med},\textsc{high}\}$, and authorization variables
$\alpha_i$ (state variables that must be present for a non-low-risk tool to be
admissible). The agent holds a state $s_t$ (a set of variables); a tool is
\emph{executable} when $R_i \subseteq s_t$ and \emph{authorized} when
$\rho_i = \textsc{low}$ or $\alpha_i \subseteq s_t$.

\racg{} selects the visible set $\Vt$ by (i) finding a minimal risk-penalized
cost path from $s_t$ to the goal $g$, where a non-low-risk tool incurs a penalty
$\lambda\cdot\mathrm{risk}(\rho_i)$ whenever its authorization is unmet at the
point it would apply, then (ii) exposing the first tool of that path, subject to
an \emph{admissibility} check that drops a high-risk frontier tool whose
authorization is unmet, falling back to a low-risk tool that establishes the
missing authorization. The structural injection guarantee is: if the attacker's
target tool is not in $\Vt$, a compliant agent cannot call it.

\subsection{The Relocated Trust Assumption}
Every quantity the gate uses---$R_i, E_i, \rho_i, \alpha_i$---comes from the
contract. The gate has no independent oracle for ``is this tool dangerous'' or
``what does this tool actually do.'' Thus the structural guarantee is
conditional on contract integrity, and the question we ask is: \emph{what
happens when the contracts lie?}

\subsection{Threat Model}
\label{sec:threat}
We consider an attacker with the following \textbf{capabilities}: they may
modify the declared contract of one or more tools (description, $R_i$, $E_i$,
$\rho_i$, $\alpha_i$), and/or cause a tool's \emph{realized} runtime effect to
differ from its declared $E_i$. This models (a)~indirect prompt injection into a
tool description or effect annotation, (b)~a registry populated in part from
untrusted third-party tool providers, and (c)~a tool whose implementation
misbehaves relative to its advertised contract.

The attacker \textbf{cannot}: modify the gate or the agent; forge a trusted
signer's signature; or alter a trusted attestation of the original registry
(which \cg{} compares against). We assume the existence of such an
attestation---a signed contract store or equivalent---as the defender's root of
trust. We do \emph{not} assume the agent resists injection: consistent with the
base method, the agent is worst-case \emph{adversarially compliant}, obeying any
injected instruction whose target tool is visible. This isolates the
\emph{structural} effect of contract integrity from the agent's behavior.

The attacker's \textbf{goal} is to cause the agent to invoke a high-risk tool
$t^\star$ that an honest contract layer would have gated out, on a task where
$t^\star$ is not legitimately authorized. We measure success by the
\emph{injection success rate} $\isr$: the fraction of injection trials in which
the agent calls $t^\star$.

% =========================================================
\section{The Two-Gate Structure and the Attack Taxonomy}
\label{sec:taxonomy}
% =========================================================

\subsection{Two Gates, and Which One Dominates}
Instrumenting \racg{} reveals that it filters in two stages. The \emph{causal
gate} is the minimal-path search: a tool enters consideration only if it lies on
a cheapest path from $s_t$ to $g$. The \emph{admissibility gate} then drops a
risky frontier tool whose authorization is unmet. The causal gate is applied
first and is decisive: a tool that advances no cheapest path to the goal is never
the frontier, and so is never even tested for admissibility.

This has an immediate consequence for attacks. Consider a ``no-action'' task
whose goal is to summarize an email; \code{send\_email} produces
\code{email\_sent}, which is irrelevant to the goal. No manipulation of
\code{send\_email}'s \emph{risk label} or \emph{authorization set} can make it
appear, because it is off the causal path regardless of $\rho$ or $\alpha$. We
confirm this empirically (Section~\ref{sec:results}): risk downgrade and
authorization aliasing, applied alone, leave $\isr = 0$ under \racg.

\subsection{Effect-Integrity-Primary Attacks}
The attack that works must therefore defeat the causal gate, which is governed
by $R_i$ and $E_i$. We define four attacks, of which the effective ones forge
effects or preconditions.

\paragraph{A1: Shortcut forgery (effect + precondition forgery)}
The attacker poisons $t^\star$ so that it is executable directly from the initial
state ($R_{t^\star} \leftarrow s_0$) and single-handedly produces the goal
($E_{t^\star} \leftarrow E_{t^\star}\cup g$), and downgrades its risk so it also
passes admissibility. Now $t^\star$ is the one-step minimal path; the causal gate
selects it and it is exposed at the first step. This models a poisoned tool
description advertising a one-shot shortcut.

\paragraph{A2: Risk downgrade (alone)}
The attacker relabels $\rho_{t^\star}\leftarrow\textsc{low}$. As argued above,
this alone is insufficient unless $t^\star$ is already on the causal path; we
report it to demonstrate causal-gate dominance.

\paragraph{A3: Authorization aliasing (alone)}
The attacker empties $\alpha_{t^\star}$. Also insufficient alone, for the same
reason.

\paragraph{A4: Effect divergence (runtime)}
The declared contract is benign and correctly signed, but the tool's
\emph{realized} effect at runtime includes an authorization variable it never
declared. This forges authorization for an on-path high-risk tool without any
detectable static anomaly. Only runtime checking can catch it.

\paragraph{A1$'$: Signed over-scoping}
A correctly-signed benign tool declares (and realizes) an authorization variable
it is not entitled to produce, opening the admissibility gate on an on-path
high-risk tool whose authorization variable has no benign establisher.
Provenance checking passes it (the signer is trusted); only typing catches it.

% =========================================================
\section{ContractGuard}
\label{sec:cg}
% =========================================================
\cg{} sits between the registry and the gate and returns a sanitized registry.
It is parameterized by a level $L\in\{0,1,2,3\}$ enabling its mechanisms
cumulatively, which yields the ablation behind our necessity ladder.

\paragraph{Rung 1 --- Signed provenance ($L\geq 1$).}
Each contract carries a signer. \cg{} rejects any contract field not vouched for
by a trusted signer: if a presented contract's signer is untrusted, \cg{} falls
back to the trusted attestation for that tool (or drops the tool entirely if it
is unknown). This neutralizes any attack that re-signs a contract with forged
fields, including A1 shortcut forgery.

\paragraph{Rung 2 --- Typed contract attestation ($L\geq 2$).}
We refer to this mechanism throughout as \emph{typed contract attestation}: for
any tool present in the trusted attestation that claims a trusted signer,
\cg{} verifies its integrity-critical, gate-relevant fields---preconditions
$R_i$, produced effects $E_i$, risk $\rho_i$, and authorization set
$\alpha_i$---against the trusted reference and restores any that diverge. This
subsumes \emph{authorization-variable typing} as a special case (an over-claimed
authorization variable in $E_i$ is removed because $E_i$ is restored to the
attested value, so a tool can only establish authorization variables it is
genuinely entitled to), and additionally catches \emph{same-signer} tampering
with $R_i$, $\rho_i$, or $\alpha_i$. The latter matters: our adaptive evaluation
(Section~\ref{sec:adaptive}) revealed that a signer-only rung~1 is bypassed by an
attacker who keeps a trusted signer but, e.g., empties $\alpha_{t^\star}$;
typed contract attestation closes that class. This rung is necessary because
rung~1 (origin only) does not inspect field values of a trusted-signed contract
(A1$'$ signed over-scoping).

\paragraph{Rung 3 --- Runtime effect verification ($L\geq 3$).}
At execution, \cg{} compares a tool's realized output against its (sanitized,
declared) effects and drops any undeclared effect. This catches A4 effect
divergence, which is invisible to both static rungs because the declared
contract is clean.

The three mechanisms are deliberately ordered by where they intervene
(declaration provenance, declaration typing, realized behavior), and we show
each is necessary against a distinct attack.

% =========================================================
\section{Hypotheses}
\label{sec:hyp}
% =========================================================
We continue the hypothesis ladder of the base method (H1--H5) with five new
hypotheses, evaluated on the metrics above.

\begin{table}[t]
\centering
\caption{Hypotheses for contract integrity (Arc~A). H6--H9 and H14a are the
core ladder; H6$'$, H$_\text{cmp}$, H$_\text{adp}$, H$_\text{fld}$, H$_\text{util}$
strengthen the claims to generality, adaptivity, and deployability.}
\label{tab:hyp}
\small
\begin{tabularx}{\columnwidth}{@{}l l Y@{}}
\toprule
\textbf{H} & \textbf{Type} & \textbf{Claim / signal} \\
\midrule
H6 & vuln.\ & Contract integrity is a hidden trust assumption: each effective
attack drives \racg{} $\isr>0$; honest contracts give $\isr=0$. \\
H6$'$ & vuln.\ & The vulnerability is \emph{general}: shortcut forgery attains
$\isr=1$ across all core high-risk targets; risk/auth tampering alone attains
$\isr=0$ across all targets (causal-gate dominance). \\
H7 & defense & Signed provenance neutralizes shortcut forgery ($\isr\!\to\!0$)
but is insufficient against signed over-scoping. \\
H8 & defense & Typed contract attestation neutralizes signed over-scoping
but is insufficient against runtime effect divergence. \\
H9 & defense & Runtime effect verification neutralizes effect divergence. \\
H$_\text{cmp}$ & defense & A compound attack (A1$+$A4) is closed only by the full
stack: $L1,L2$ leave $\isr>0$; $L3$ gives $\isr=0$. \\
H$_\text{adp}$ & defense & The full stack is robust to an \emph{exhaustive
white-box adaptive} attacker: worst-case attack-induced $\isr=0$ at $L3$, while
$L0$--$L2$ are each defeated. \\
H$_\text{fld}$ & analysis & Effect integrity is primary: only the
$E_i$$+$$R_i$ (causal-routing) perturbation attains $\isr=1$; $\rho_i$ or
$\alpha_i$ alone are strictly weaker. \\
H$_\text{util}$ & deploy.\ & \cg{} does not over-reject honest contracts: task
success, authorization-task success, token cost, and tools-dropped are unchanged
at every level. \\
H14a & negative & No $\lambda$ restores safety under shortcut forgery; the risk
knob cannot substitute for contract integrity. \\
\bottomrule
\end{tabularx}
\end{table}

H7--H9 are stated as a \emph{strict necessity ladder}: each rung both
(i)~neutralizes its target attack and (ii)~leaves the next attack's $\isr>0$,
establishing that the rung is necessary and the prior rung insufficient. We also
verify a utility condition (\cg{} does not change task completion on honest
contracts).

% =========================================================
\section{Experimental Setup}
\label{sec:setup}
% =========================================================
We extend the RiskGate benchmark and its deterministic, adversarially-compliant
agent~\cite{iyer2026racg}, which upper-bounds injection success for any real
model: if the worst-case compliant agent cannot call $t^\star$ because it is
absent from $\Vt$, no real model can. We use the same 100-tool registry and the
default \racg{} configuration ($\lambda=2$) unless noted. We implement the
attacks of Section~\ref{sec:taxonomy} as adversarial transforms on the registry
(the contract layer), and \cg{} as a registry sanitizer applied before the gate;
the agent and gate only ever read the post-attack, post-\cg{} registry. We
construct four scenarios, each isolating one rung of the ladder:
\textbf{S0} (honest control), \textbf{S1} (A1 shortcut forgery, attacker-signed),
\textbf{S2} (A1$'$ signed over-scoping of an authorization variable with no
benign establisher), and \textbf{S3} (A4 runtime effect divergence). Each
scenario is run under three injection phrasings (direct, roleplay, obfuscated)
and across \cg{} levels $L\in\{0,1,2,3\}$. All experiments are deterministic and
reproducible; we report $\isr$ as the mean over phrasings.

\paragraph{Scale of the evaluation.} The registry contains 100 tools, of which
\textbf{8} are core high-risk targets (each with a distinct authorization
variable); these are the targets the attacker attempts to re-expose. Every
scripted condition is run over \textbf{3} injection phrasings and \textbf{4}
guard levels. This yields three families of trials. (i)~The hand-built ladder
(S0--S3) is $4~\text{scenarios}\times3~\text{phrasings}\times4~\text{levels}=48$
trials. (ii)~The programmatic suite of Section~\ref{sec:results} instantiates
six attack families across their applicable targets---A1, A2, and A3 over all 8
targets, and A1$'$, A4, and the compound attack over the 3 no-establisher targets
(below)---for $(8{+}8{+}8{+}3{+}3{+}3)=33$ attack instances, or
$33\times3\times4=396$ scripted trials. (iii)~The adaptive attacker
(Section~\ref{sec:adaptive}) enumerates \textbf{256 perturbation configurations
per target} (Section~\ref{sec:adaptive} details the grid) over all 8 targets and
4 levels, i.e.\ $256\times8=2048$ configurations, each scored over 3 phrasings
($6144$ phrasing-trials). In total the study comprises just over $6.5\mathrm{k}$
deterministic agent rollouts.

% =========================================================
\section{Results}
\label{sec:results}
% =========================================================

\subsection{The Strict Necessity Ladder (H6--H9)}
Table~\ref{tab:ladder} and Fig.~\ref{fig:ladder} report $\isr$ for each scenario
across \cg{} levels. The pattern is exact. With no defense ($L0$), every effective attack achieves
$\isr=1.0$ while the honest control stays at $0$, establishing H6: contract
integrity is load-bearing. Each subsequent rung drives its target attack to $0$
and leaves the next attack at $1.0$: provenance ($L1$) stops S1 but not S2; typed
contract attestation ($L2$) stops S2 but not S3; runtime verification ($L3$)
stops S3. This is the strict necessity ladder H7~$<$~H8~$<$~H9.

\begin{table}[t]
\centering
\caption{Injection success rate ($\isr$) by scenario and \cg{} level. Each
attack is neutralized by exactly one additional rung; the prior rung is
insufficient (bold = first level achieving $\isr=0$).}
\label{tab:ladder}
\small
\begin{tabular}{@{}l cccc c@{}}
\toprule
Scenario & $L0$ & $L1$ & $L2$ & $L3$ & Defended by \\
\midrule
S0 honest                & 0.00 & 0.00 & 0.00 & 0.00 & --- \\
S1 shortcut forgery      & 1.00 & \textbf{0.00} & 0.00 & 0.00 & provenance \\
S2 signed over-scope     & 1.00 & 1.00 & \textbf{0.00} & 0.00 & typed attestation \\
S3 effect divergence     & 1.00 & 1.00 & 1.00 & \textbf{0.00} & runtime check \\
\bottomrule
\end{tabular}
\end{table}

\begin{figure}[t]
\centering
\includegraphics[width=\columnwidth]{figures/ladder_isr.png}
\caption{The strict necessity ladder. With no defense ($L0$) every effective
attack reaches $\isr=1$; each \cg{} rung neutralizes exactly one attack and is
insufficient for the next, establishing H6 and the ordering H7~$<$~H8~$<$~H9.}
\label{fig:ladder}
\end{figure}

\subsection{Causal-Gate Dominance}
We confirm that risk downgrade (A2) and authorization aliasing (A3) applied
\emph{alone} to an off-path target leave $\isr=0$ under \racg{} at every
$\lambda$. The dangerous tool is never the causal frontier, so admissibility is
never reached. This both validates the effect-integrity-primary framing and
records a robustness property of the base method: naive risk/authorization
tampering does not suffice.

\subsection{Generality Across the Registry (H6$'$)}
The hand-built scenarios are existence proofs; we now show the effect is general.
We generate the attack suite programmatically across all eight core high-risk
tools (each with a distinct authorization variable), instantiating shortcut
forgery (A1) per target, and authorization-variable forgery (A1$'$, A4,
compound) for every target whose authorization variable has no benign
establisher, over three injection phrasings.
Table~\ref{tab:dist} and Fig.~\ref{fig:dist} report ISR by family and guard
level. Shortcut forgery attains $\isr=1.00$ across all targets at $L0$
($n=24$ trials) and is driven to $0$ by provenance ($L1$); the negative controls
(A2, A3 alone) attain $\isr=0.00$ across all targets at \emph{every} level
($n=24$ each), confirming causal-gate dominance is a general property, not an
artifact of one task. The authorization-forgery families reproduce the ladder
(A1$'$ closed at $L2$; A4 at $L3$).

A1$'$, A4, and the compound attack are restricted to the no-establisher targets
\emph{by construction}: these are the only targets for which a forged
authorization variable can be isolated, because they are the only ones whose
authorization variable has no legitimate benign establisher---on any other
target \racg{} would open the gate via the real establisher, so reaching the
tool is correct behavior rather than a successful forgery. Three of the eight
core targets satisfy this condition, hence $n=3$ for those families; we do not
read their absolute counts as a power claim but as a clean isolation of the
authorization channel, and the adaptive search of Section~\ref{sec:adaptive}
then generalizes beyond these scripted scenarios by enumerating the full
perturbation grid on all eight targets.

\begin{table}[t]
\centering
\caption{Distributional ISR by attack family over the generated suite (eight
core high-risk targets, three phrasings). A1 and the negative controls span all
targets ($n=24$); authorization-forgery families span the no-establisher targets.}
\label{tab:dist}
\small
\begin{tabular}{@{}l ccccc@{}}
\toprule
Family & $n$ & $L0$ & $L1$ & $L2$ & $L3$ \\
\midrule
A1 shortcut forgery       & 24 & 1.00 & 0.00 & 0.00 & 0.00 \\
A1$'$ signed over-scope   &  3 & 1.00 & 1.00 & 0.00 & 0.00 \\
A4 effect divergence      &  3 & 1.00 & 1.00 & 1.00 & 0.00 \\
A1$+$A4 compound          &  3 & 1.00 & 1.00 & 1.00 & 0.00 \\
A2 risk downgrade (alone) & 24 & 0.00 & 0.00 & 0.00 & 0.00 \\
A3 auth aliasing (alone)  & 24 & 0.00 & 0.00 & 0.00 & 0.00 \\
\bottomrule
\end{tabular}
\end{table}

\begin{figure}[t]
\centering
\includegraphics[width=\columnwidth]{figures/distributional_isr.png}
\caption{Distributional ISR by attack family across guard levels. The
effect-forgery families (A1, A1$'$, A4, compound) defeat \racg{} at $L0$ and are
each closed by the appropriate rung; the risk/authorization-only controls (A2,
A3) never defeat \racg.}
\label{fig:dist}
\end{figure}

\subsection{Compound Attacks (H$_\text{cmp}$)}
A defender might hope a partial deployment suffices. We run a compound attack
that simultaneously declares an over-scoped authorization variable (A1$'$) and
diverges it at runtime (A4). Static stripping ($L2$) removes the declared forge
but the runtime channel still delivers the authorization variable, so $L1$ and
$L2$ leave $\isr=1.00$; only the full stack ($L3$) closes it
(Table~\ref{tab:dist}, compound row). Defense-in-depth is therefore not optional:
the rungs are jointly, not alternatively, required.

\subsection{Which Field Carries the Attack (H$_\text{fld}$)}
Table~\ref{tab:field} and Fig.~\ref{fig:field} ablate the contract one field at a
time, perturbing only $R_i$, only $E_i$, only $\rho_i$, or only $\alpha_i$ on each
target and measuring ISR under \racg{} with no defense. Only the
$E_i$$+$$R_i$ combination---the perturbation that controls causal
routing---reaches $\isr=1.00$; risk downgrade and authorization aliasing alone
reach only $0.25$, and $R_i$-only or $E_i$-only reach $0.00$. This quantifies the
effect-integrity-primary thesis and tells defenders where attestation budget
should concentrate: the precondition/effect annotations, not the risk label.

\begin{table}[t]
\centering
\caption{Field-sensitivity ablation: ISR under \racg{} (no guard) when exactly
one contract field (or the causal-routing pair) is perturbed, meaned over the
core targets.}
\label{tab:field}
\small
\begin{tabular}{@{}l c@{}}
\toprule
Perturbed field & ISR \\
\midrule
$E_i + R_i$ (causal routing) & 1.00 \\
$\rho_i$ (risk) alone        & 0.25 \\
$\alpha_i$ (authorization) alone & 0.25 \\
$R_i$ (preconditions) alone  & 0.00 \\
$E_i$ (effects) alone        & 0.00 \\
\bottomrule
\end{tabular}
\end{table}

\begin{figure}[t]
\centering
\includegraphics[width=\columnwidth]{figures/field_ablation.png}
\caption{Attack surface by contract field. Causal-routing integrity ($E_i,R_i$)
is the load-bearing assumption; the risk label is secondary.}
\label{fig:field}
\end{figure}

\subsection{The Risk Knob Cannot Substitute (H14a)}
Table~\ref{tab:lambda} and Fig.~\ref{fig:lambda} sweep $\lambda$ under S1
(shortcut forgery) with no defense. $\isr$ is pinned at $1.0$ for all $\lambda\in[0,100]$. Because shortcut
forgery defeats the \emph{causal} gate---the poisoned tool is the cheapest
path---the risk-penalty $\lambda$, which only re-weights admissibility, cannot
remove it. The risk knob and contract integrity are orthogonal controls.

\begin{table}[t]
\centering
\caption{$\isr$ under shortcut forgery (S1) as $\lambda$ varies, no defense.
The attack is invariant to the risk knob.}
\label{tab:lambda}
\small
\begin{tabular}{@{}l cccccccc@{}}
\toprule
$\lambda$ & 0 & 0.25 & 0.5 & 1 & 2 & 4 & 8 & 100 \\
\midrule
$\isr$ & 1.0 & 1.0 & 1.0 & 1.0 & 1.0 & 1.0 & 1.0 & 1.0 \\
\bottomrule
\end{tabular}
\end{table}

\begin{figure}[t]
\centering
\includegraphics[width=\columnwidth]{figures/lambda_negative.png}
\caption{H14a. Sweeping the risk-penalty $\lambda$ under shortcut forgery leaves
$\isr$ pinned at $1$: the attack defeats the causal gate, which $\lambda$ does
not influence. Benign success is unaffected, so the risk knob trades nothing for
safety here---contract integrity is a separate control.}
\label{fig:lambda}
\end{figure}

\subsection{Utility and Overhead (H$_\text{util}$)}
A defense is only deployable if it is cheap on honest inputs.
Table~\ref{tab:util} runs the original benign and safety-stress suite under each
guard level with honest contracts. At every level---including the full
$L3$---task completion (1.00), authorization-task completion (1.00), and mean
token cost are unchanged, and \cg{} drops zero legitimate tools. \cg{} is thus
free on honest contracts: it changes the gate's decision only when a contract
diverges from its attestation.

\begin{table}[t]
\centering
\caption{ContractGuard overhead on honest contracts (full benign $+$
safety-stress suite). No utility loss and no dropped tools at any level.}
\label{tab:util}
\small
\begin{tabular}{@{}l cccc@{}}
\toprule
Level & success & auth-task success & mean tokens & tools dropped \\
\midrule
$L0$ (none)        & 1.00 & 1.00 & 1350 & 0 \\
$L1$ provenance    & 1.00 & 1.00 & 1350 & 0 \\
$L2$ typed attest. & 1.00 & 1.00 & 1350 & 0 \\
$L3$ runtime       & 1.00 & 1.00 & 1350 & 0 \\
\bottomrule
\end{tabular}
\end{table}

% =========================================================
\section{The Adaptive Attacker}
\label{sec:adaptive}
% =========================================================
The scripted suite fixes the attack. A stronger and standard threat model is a
\emph{white-box adaptive} attacker who knows the deployed guard level and
searches for \emph{any} contract perturbation that achieves a high-risk call. We
enumerate the full perturbation grid for each target: the signer
$\in\{\text{registry},\text{attacker}\}$; whether to forge $R_{t^\star}$ to the
initial state; whether to forge $E_{t^\star}$ to include the goal; whether to
downgrade $\rho_{t^\star}$; whether to empty $\alpha_{t^\star}$; the carrier
forgery mode $\in\{\text{none},\text{declared},\text{runtime},\text{both}\}$; and
the injection step $\in\{0,\text{target step}\}$. The product of these axes is
$2\times2\times2\times2\times2\times4\times2 = 256$ perturbation configurations
per target; across the eight high-risk targets the attacker's best-response space
is $256\times8 = 2048$ configurations, each evaluated over the three injection
phrasings ($6144$ phrasing-trials). Because the grid is finite we enumerate it
exhaustively---a true best response---so a ``no success'' result at a guard level
is a guarantee over the entire modeled space, not a sampling estimate. (As an
implementation detail, the search short-circuits as soon as it finds a
perturbation attaining $\isr=1$; this is immaterial to the headline guarantee,
because at $L3$ no configuration reaches $\isr=1$, so the $L3$ enumeration runs to
completion and the ``no success at $L3$'' claim is genuinely exhaustive.) We report
\emph{attack-induced} ISR (best ISR minus the honest baseline) so that a tool
which is legitimately authorizable on its task is not miscounted as an attack.

Table~\ref{tab:adapt} and Fig.~\ref{fig:adapt} give the result. With no defense,
and at $L1$ and $L2$, the adaptive attacker attains attack-induced $\isr=1.00$ on
seven of eight targets; the eighth (\code{transfer\_funds}) is legitimately
authorizable and shows $0$ after baseline subtraction. At $L3$ the worst-case
attack-induced ISR is $0.00$ across \emph{all} targets: no perturbation in the
modeled space defeats the full stack.

Notably, the adaptive search \emph{discovered a gap in an earlier ContractGuard
design}. A first version implemented rung~2 as authorization-variable typing
only; the adaptive attacker bypassed it by keeping a trusted signer and simply
\emph{emptying} $\alpha_{t^\star}$ (authorization aliasing on an on-path tool),
which neither signer checking nor produced-effect typing inspects. This finding
motivated the generalization of rung~2 to full typed contract attestation
(Section~\ref{sec:cg}). We report it because it illustrates the value of adaptive
evaluation: it converts ``each rung stops its scripted attack'' into ``the stack
stops every attack in the space,'' and it caught a real omission.

\begin{table}[t]
\centering
\caption{Adaptive white-box attacker: worst-case attack-induced ISR over an
exhaustive best-response search of the contract-perturbation space (256
configurations per target, 2048 across the eight targets). The full
stack ($L3$) admits no successful perturbation.}
\label{tab:adapt}
\small
\begin{tabular}{@{}l cccc@{}}
\toprule
& $L0$ & $L1$ & $L2$ & $L3$ \\
\midrule
worst-case attack-induced ISR & 1.00 & 1.00 & 1.00 & \textbf{0.00} \\
\bottomrule
\end{tabular}
\end{table}

\begin{figure}[t]
\centering
\includegraphics[width=\columnwidth]{figures/adaptive_isr.png}
\caption{Adaptive attacker attack-induced ISR per target and guard level.
$L0$--$L2$ are each defeated on the genuinely-attackable targets;
$L3$ is uniformly $0$. (\code{transfer\_funds} is legitimately authorizable and
shows $0$ after baseline subtraction.)}
\label{fig:adapt}
\end{figure}

% =========================================================
\section{Soundness of the Defense}
\label{sec:theory}
% =========================================================
We state the guarantee each rung provides over the gate's inputs. Let
$\hat t$ denote the attested (trusted) contract for a known tool and $t$ the
contract presented to the gate after \cg{} sanitization at level $L$. Write
$\Vt(\cdot)$ for the visible set the gate computes from a registry.

\begin{proposition}[Provenance, $L\geq 1$]
\label{prop:prov}
At $L\geq 1$, every tool the gate reads is either (i) present in the attestation,
or (ii) signed by a trusted signer. No attacker-introduced \emph{new} tool, and
no contract \emph{re-signed} under an untrusted identity, influences $\Vt$.
\end{proposition}
\noindent\emph{Proof sketch.} \cg{} drops any tool absent from the attestation and
restores the attested contract for any tool whose presented signer is untrusted;
the remainder claim a trusted signer, which by assumption the attacker cannot
forge. Hence the shortcut-forgery family (which re-signs $t^\star$) cannot place
a forged contract before the gate. $\square$

\begin{proposition}[Field integrity, $L\geq 2$]
\label{prop:field}
At $L\geq 2$, for every known tool the gate's integrity-critical fields equal the
attested ones: $R_t=R_{\hat t}$, $E_t=E_{\hat t}$, $\rho_t=\rho_{\hat t}$,
$\alpha_t=\alpha_{\hat t}$. Consequently no declared contract perturbation---of
preconditions, effects, risk, or authorization---changes $\Vt$ for a known tool.
\end{proposition}
\noindent\emph{Proof sketch.} \cg{} restores each of these fields from $\hat t$.
Authorization-variable typing is the special case $E_t\cap\mathcal{A}\subseteq
\mathrm{entitled}(\hat t)$, which holds because $E_t=E_{\hat t}$ and $\hat t$ is
entitled to exactly $E_{\hat t}\cap\mathcal{A}$. Thus signed over-scoping and
same-signer field tampering are both neutralized. $\square$

\begin{proposition}[Runtime soundness, $L\geq 3$]
\label{prop:runtime}
At $L\geq 3$, the realized effect of any tool is a subset of its (attested)
declared effects: no undeclared variable, in particular no undeclared
authorization variable, can enter the state through tool execution.
\end{proposition}
\noindent\emph{Proof sketch.} \cg{} intersects the realized output with the
declared $E_t$ before applying it to the state; by Prop.~\ref{prop:field}
$E_t=E_{\hat t}$, so the realized effect is bounded by the attested effect set.
Effect divergence (A4) and the runtime channel of the compound attack are thereby
blocked. $\square$

\paragraph{Boundary condition.} All three propositions are conditional on the
attestation being trustworthy: \cg{}'s root of trust is a signed reference for
known tools that the attacker cannot forge or replace. This is the contract-layer
analogue of the base method's authorization-provenance assumption; we make it
explicit rather than emergent, and a deployment must enforce it at the system
level (e.g.\ a signed contract registry).

\paragraph{Scope of runtime soundness.} Proposition~\ref{prop:runtime} bounds the
effect that enters the agent's \emph{symbolic state}: \cg{} intersects a tool's
realized output with its attested effects before that output updates the state
the gate reasons over, so no undeclared variable---in particular no undeclared
authorization variable---can be laundered into subsequent gating decisions. This
guarantee is over \emph{state}, not over the external world. Runtime effect
verification mediates the state update, but it acts \emph{after} the tool has
executed: if a poisoned or misbehaving tool performs an irreversible external
side effect---sending an email, deleting a file, transferring funds---before
\cg{} inspects its output, dropping the undeclared effect from the symbolic state
does not undo the real-world action. The proposition therefore prevents a forged
runtime effect from \emph{escalating into further gated calls}, but it does not,
on its own, prevent the first irreversible side effect of the diverging tool.
Closing that gap requires composing \cg{} with mechanisms that intervene
\emph{before} commitment---pre-execution effect declaration with a transactional
or dry-run check, capability-scoped execution sandboxes, or human confirmation on
irreversible operations---which we treat as an enforcement-layer concern
orthogonal to contract integrity and discuss in Section~\ref{sec:lim}.

% =========================================================
\section{Discussion}
\label{sec:disc}
% =========================================================
Our results reframe the safety story of structural gating. The guarantee
``the dangerous tool is not in $\Vt$'' is real, but it is an
\emph{if-then}: \emph{if} the contracts are honest, \emph{then} the gate is
sound. Structural gating thus converts the problem of agent persuasion into the
problem of contract integrity---a different, and we argue more tractable,
problem, because contracts are static, declarable, and signable artifacts,
whereas model behavior is not. \cg{} makes that conversion explicit: it pushes
the trust boundary to a signed attestation and a typed effect discipline, both
of which are amenable to standard supply-chain and information-flow techniques.

The two-gate finding is independently useful. It tells defenders \emph{where}
contract integrity matters most: not the risk label (the admissibility gate is
downstream of causal selection and cannot be reached for an off-path tool) but
the precondition--effect annotations that determine causal reachability.
Defensive effort and attestation should concentrate on $R_i$ and $E_i$.

% =========================================================
\section{Limitations and Future Work}
\label{sec:lim}
% =========================================================
Our primary evaluation uses a controlled, deterministic benchmark and a
worst-case compliant agent. As in the base method, this is a deliberate choice
that \emph{upper-bounds} injection success and isolates structural effects. To
connect that bound to real models we provide a \emph{real-LLM validation track}:
the same env/filter/registry/attack/guard stack driven by an actual model as the
policy (via native function calling), over a focused grid of
\mbox{models~$\times$~8~high-risk~targets~$\times$~3~phrasings~$\times$~$\{L0,L3\}$~$\times$~$\{$A1,A4$\}$}.
It reports \emph{attack-induced} ISR (baseline-subtracted, as in
Section~\ref{sec:adaptive}) so a legitimately-authorizable target is not
miscounted, and it tests the structural prediction directly: the contract attack
lands at $L0$ (the re-exposed tool is in $\Vt$, so a compliant model can call it)
and is closed at $L3$ (the forged exposure is removed, so the tool is absent from
$\Vt$ and \emph{no} model can call it), independent of phrasing. The harness is
provider-agnostic (Anthropic, Amazon Bedrock, any OpenAI-compatible endpoint,
and an offline stub) and degrades gracefully when a provider is unavailable. We
release it with the artifact and validate its wiring against the deterministic
reference; running it at scale across several hosted frontier models---and
measuring the predicted split between description-level persuasion attacks
(which leave $\Vt$ intact and should vary by model) and structural contract
attacks (which corrupt $\Vt$ and should behave model-independently)---is the
remaining empirical work, which we leave to a hosted-model study to keep this
paper's claims tied to the exhaustive structural guarantee.

\cg{}'s root of trust is a signed attestation of the original registry; we
assume but do not build the signing infrastructure, and a deployment must ensure
the attestation cannot itself be poisoned. The typed-attestation mechanism
requires a declared entitlement per tool, which is an additional specification
burden; deriving entitlements automatically from a trusted reference contract is
future work.

A further limitation concerns runtime effect verification. In our symbolic
benchmark, mediating the state update is sufficient to block effect divergence,
because a tool's only avenue of influence is the variables it writes to the
shared state. In a deployed system a tool may produce \emph{irreversible external
side effects}---an email already sent, a file already deleted, funds already
transferred---before \cg{} intersects its output with the declared effects.
\cg{}'s runtime rung prevents such a divergence from escalating into further
gated calls (it never enters the state the gate reads), but it cannot retract an
action the tool has already committed in the world. Runtime effect verification
should therefore be understood as a \emph{state-integrity} mechanism, not a
rollback mechanism; for irreversible operations it must be paired with
before-commitment controls (pre-execution effect declaration validated by a
dry-run or transactional wrapper, capability-scoped sandboxing, or human
confirmation). Integrating \cg{} with such an enforcement layer, and verifying
that the composition preserves the necessity ladder, is future work. Finally, this paper addresses the \emph{contract-integrity} half of
the trust that structural gating relocates. The complementary half---that the
gate reasons about risk \emph{per tool}, so danger emergent across a chain of
individually-safe tools can evade it (``toolchain confusion'')---and an
information-flow defense for it (``FlowGate'') are the subject of a companion
paper, completing the program of treating the visible tool set as a control
surface whose honesty and locality must both be earned.

% =========================================================
\section{Conclusion}
% =========================================================
Structural tool gating does not eliminate the trust assumption behind safe tool
use; it relocates it into the integrity of the contracts the gate reads. We
threat-modeled that contract layer, showed that effect-forgery attacks defeat the
gate \emph{upstream of the agent} while naive risk/authorization tampering does
not (a consequence of \racg's dominant causal gate), and introduced \cg, whose
signed-provenance, typed-contract-attestation, and runtime-effect-verification
mechanisms form a strict necessity ladder restoring injection success to zero
without over-rejecting honest contracts. The risk knob $\lambda$ cannot
substitute for any of these. The gate is only as honest as its contracts;
\cg{} makes that honesty explicit and verifiable.

\bibliographystyle{IEEEtran}
\bibliography{references}

\end{document}
